\section{Problem statement and objectives}\label{sec:problem}

The structure of the model is fixed in advance by the robot's kinematics and its sensor geometry, so there is no
structure to learn. What is left is the inference: the belief cannot be computed in closed form, so it is
approximated by a set of samples, and the estimate is only as good as that set. How that set is proposed, weighted and resampled is the
first thing measured here.

The motion model and the measurement model each carry a noise level, and in practice these are almost always set
by hand. In this project I estimate them instead, from the data the robot itself produced. That is the second
thing measured here, and the weights the filter already computes are what make it possible.

\paragraph{Sampling.} A filter must choose a proposal distribution, a rule for when to resample, a resampling
scheme and a number of particles, and these are usually inherited from whatever implementation is at hand. I
built a testbed in which they are interchangeable parts of one filter loop, so that one can be varied while the
world, the trajectory and the random seed stay fixed. The filters themselves are existing open-source
implementations: the localization filters and resampling schemes of Elfring et al.\ \cite{elfring2021particle},
and FastSLAM 1.0 from PythonRobotics \cite{sakai2018pythonrobotics}.

Two of the four are the subject of this report:
\begin{itemize}
  \item \emph{the resampling scheme}: multinomial, residual, stratified and systematic resampling
        \cite{douc2005comparison,hol2006resampling}, which differ in the variance they add when the set is
        redrawn;
  \item \emph{when to resample}: at every step; never; when the effective sample size falls below a threshold
        (0.2, 0.5 or 0.8 of $N$); or when the largest weight rises above one (0.1, 0.2 or 0.5).
\end{itemize}
The other two are held fixed so that they cannot confound the comparison. The proposal is always the motion
model, which draws from the prior, and the number of particles is fixed within each comparison.

\paragraph{Parameter learning.} The noise values are parameters of the network's two conditional distributions,
and the weights the filter already computes estimate the marginal likelihood of the observations under them,
so they can be estimated from a recorded run instead of assumed. This asks not how well
the filter infers under a given model, but which parts of the model the data can determine at all. What is
tested:
\begin{itemize}
  \item whether the four noise parameters can be recovered by maximizing \eqref{eq:mle}, and how the recovered
        values compare with the values the world used;
  \item which of them the data identify: whether the likelihood has a maximum, and whether independent recordings
        agree on where it is;
  \item what the inference must be capable of first: how the estimate depends on the number of particles, on the
        resampling scheme, and on whether the filter tracks or localizes from scratch.
\end{itemize}
